\section{Label aggregation and misspecified model}
\label{sec:mis-spec}

The logistic link provides clean interpretation and results, but we can
move beyond it to more realistic cases where labelers use
distinct links; to allow precise statements, we still assume the
same linear term $x \mapsto \<\theta\opt, x\>$ for each labeler's
generalized linear model. We study generalizations of the maximum
likelihood and majority vote estimators~\eqref{eq:maximum-likelihood}
and~\eqref{eq:logistic-majority}, highlighting dependence on
link fidelity.
In this setting, there are $m$ (unknown and possibly distinct)
link functions $\sigma_i\opt$, $i=1,2, \ldots, m$.
We show that the majority-vote estimator $\mve$ enjoys better robustness to model mis-specification than the non-aggregated estimator $\tmle$,
though both use identical losses.
In particular, our main result in this section implies
%\begin{equation*}
%\sqrt{n}(\what{u}\lr_{n,m} - u\opt) \cd
%\normal\left(0, c\lr \Sigma \cdot (1 + o_m(1))\right)
%~~ \mbox{while} ~~
%\sqrt{n}(\what{u}\mv_{n,m} - u\opt) \cd
%\normal\left(0, m^{-1/2} c\mv \Sigma \cdot (1 + o_m(1))\right),
%\end{equation*}
\begin{equation*}
	\sqrt{n}(\what{u}\lr_{n,m} - u\opt) \cd
	\normal\left(0, c\lr \Sigma \right)
	~~ \mbox{while} ~~
	\sqrt{n}(\what{u}\mv_{n,m} - u\opt) \cd
	\normal\left(0, m^{-1/2} c\mv \Sigma \right),
\end{equation*}
where $c\lr$ and $c\mv$ are problem-dependent constants. %% that depend only on the links $\sigma$,
%% $t\opt = \ltwos{\theta\opt}$, and $\Sigma = I - u\opt {u\opt}^\top$ when $X \sim
%% \normal(0, I_d)$.  
In contrast to the previous section, the majority-vote
estimator enjoys roughly $\sqrt{m}$-faster rates than the non-aggregated
estimator, maintaining its (slow) improvement with $m$, which the
MLE loses to misspecification.

To set the stage for our results, we define the
general link-based loss
\begin{equation*}
  \loss_{\sigma,\theta} (y \mid x) \defeq -\int_0^{y \<\theta, x\>} \sigma(-v) dv.
\end{equation*}
We then consider the general multi-label and majority-vote
estimators using the loss $\loss_{\sigma,\theta}$,
\begin{equation}
  \label{eqn:general-estimators}
  \gme(\sigma) \defeq \argmin_\theta \Prb_{n,m} \loss_{\sigma, \theta},
  \qquad
  \mve(\sigma) \defeq \argmin_\theta \wb{\Prb}_{n,m} \loss_{\sigma, \theta}.
\end{equation}
For the logistic link $\sigma = \sigma\lr$, we recover the
results for the 
logistic loss
$\loss_{\sigma\lr,\theta} (y \mid x) = \loss_\theta\lr(y \mid x)$ in Section~\ref{sec:well-sepcified-mod}.
We suppress the link dependence
to write $\gme, \mve$ when the context is clear.
%\begin{align*}
%	\loss_{\sigma\lr,\theta} (y \mid x)  = -\int_0^{y \<\theta, x\>} \frac{1}{1 +e^{v}} dv = \log(1 + e^{-y\<x, \theta\>}) = \loss_\theta\lr(y \mid x).
%\end{align*}

%When $\sigma = \sigma\lr$, we recover the same estimators as in~\Cref{sec:well-sepcified-mod}, that is, we have $\gme(\sigma) = \tmle$ and $\mve(\sigma\lr) = \mve$. For the majority vote estimator, we will suppress the dependence on the link $\sigma$ in the loss function, and just write $\mve$ when the context is clear. \ha{why just for MV estimator? might confuse readers.} Our goal is to understand the performance of $	\gme(\sigma)$ and $ \mve$.

\subsection{Master results} \label{sec:master-results}

%%To characterize the behavior of multiple label estimators versus majority
%% vote, 
We provide master results as a foundation for our convergence rate analyses,
and by a bit of notational chicanery, we consider the majority vote and
multiple (non-aggregated) simultaneously.
%
When the estimator uses the majority vote
$\majY$, let
\begin{equation*}
  \varphi_m(t) = \rho_m(t) \ind \{t \geq 0\} + (1 - \rho_m(t)) \ind\{t < 0\},
  ~~ \mbox{where} ~~
  \rho_m(t) \defeq
  \P(\majY = \sgn(t) \mid \<X, \theta\opt\> = t),
\end{equation*}
and in the case that the estimator uses each label
from the $m$ labelers, let
\begin{equation*}
	\varphi_m(t) = \frac{1}{m} \sum_{j = 1}^m \sigma_j\opt(t)
	= \frac{1}{m} \sum_{j = 1}^m \P(Y_j = 1 \mid \<X, \theta\opt\> = t).
\end{equation*}
In either case, we then see that the population loss with
the link-based loss $\loss_{\sigma,\theta}$ becomes
\begin{equation}
	\label{eqn:generic-loss}
	L(\theta, \sigma) = \E\left[\loss_{\sigma,\theta}(1 \mid X)
	\varphi_m(\<X, \theta\opt\>)
	+ \loss_{\sigma,\theta}(-1 \mid X)
	(1 - \varphi_m(\<X, \theta\opt\>))\right],
\end{equation}
where we have taken a conditional expectation given $X$. We assume
Assumption~\ref{assumption:x-decomposition} holds, that $X$ decomposes into
the independent sum $X= Zu\opt + W$ with $W \perp u\opt$, and the true link
functions $\sigma_j\opt \in \Flink$.
%
We further impose the following
assumption for the model link, which guarantees a
non-vanishing gap to the margin when $|Z|$ is large.
\begin{assumption}
  \label{assumption:model-link}
  For $s \in \{-1, 1\}$,
  the model link function $\sigma$ satisfies
  $\lim_{t \to \infty} \sigma(s \cdot t) = \half + sc$ for a constant
  $0 < c \le \half$ and is a.e.\ differentiable.
\end{assumption}
\paragraph{Minimizer of the population loss.}
We first characterize the (unique)
minimizers $\theta\opt_L
\defeq \argmin_\theta L(\theta, \sigma)$ of the population
loss~\eqref{eqn:generic-loss}.
%
We use the familiar ansatz that $\theta\opt_L$ aligns with $u\opt =
\theta\opt / \ltwo{\theta\opt}$, so $\theta = t u\opt$.
%
Using the
formulation~\eqref{eqn:generic-loss}, we see that for $t\opt \defeq
\ltwo{\theta\opt}$ and $\theta = t u\opt$,
\begin{align}
	\nabla L(\theta, \sigma) &
	= -\E[\sigma(-\<\theta, X\>) X
	\varphi_m(\<X, \theta\opt\>)]
	+ \E[\sigma(\<\theta, X\>) X (1 - \varphi_m(\<X, \theta\opt\>))]
	\nonumber \\
	& = -\E[\sigma(- t Z) X \varphi_m(t\opt   Z)]
	+  \E[\sigma(tZ) X (1 - \varphi_m(t\opt Z))] \nonumber \\
	& = \left(-\E[\sigma(- tZ) Z \varphi_m(t\opt Z)] 
	+ \E[\sigma(tZ) Z (1 - \varphi_m(t\opt Z))] \right) u\opt
	= h_m(t) u\opt, \label{eqn:generic-gradient}
\end{align}
where the final line uses the decomposition $X = Z u\opt + W$ for the random
vector $W \perp u\opt$ independent of $Z$, and we recall
expression~\eqref{eqn:h-centering} to define the
\emph{calibration function}
\begin{equation}
  \label{eqn:h-m-func}
  h_{t\opt, m}(t) \defeq \E[\sigma(tZ) Z (1 - \varphi_m(t\opt Z))]
  - \E[\sigma(-t Z) Z \varphi_m(t\opt Z)].
\end{equation}
The function $h$ measures the ``calibration'' gap between the
hypothesized link $\sigma$ and the label probabilities
$\varphi_m$. %% , functioning the approximately ``calibrate'' $\sigma$ to the observed probabilities.
If we presume that a solution to $h_{t\opt,m}(t) = 0$ exists, then
%% evidently 
$t u\opt$ minimizes $L(\theta, \sigma)$. In fact,
such a solution exists and is unique
(see Appendix~\ref{sec:proof-generic-solutions} for a proof):

\newcommand{\Hessian}{H_L}  % Hessian matrix (as function of t)
\newcommand{\hessfunc}{\mathsf{he}}  % Hessian function thing
\newcommand{\linkerr}{\mathsf{le}}  % Link error function
\newcommand{\varfunc}{C_m}  % variance function

\begin{lemma}
  \label{lemma:generic-solutions}
  Let Assumption~\ref{assumption:x-decomposition} hold and $h = h_{t\opt,m}$
  be the gap function~\eqref{eqn:h-m-func}. Then there is a unique solution
  $t_m > 0$ to $h(t) = 0$, and
  $\theta\opt_L = t_m u\opt$ uniquely minimizes
  the loss~\eqref{eqn:generic-loss}. Define
  \begin{equation}
    \begin{split}
      \Hessian(t) &
      \defeq \E[(\sigma'(-tZ) \varphi_m(t\opt Z) + \sigma'(tZ)
	(1 - \varphi_m(t\opt Z))) Z^2 ] u\opt {u\opt}^\top \\
      & \qquad
      ~ + \E[\sigma'(-tZ) \varphi_m(t\opt Z)
	+ \sigma'(tZ) (1 - \varphi_m(t\opt Z))]
      \projperp \Sigma \projperp.
    \end{split}
    \label{eqn:parameterized-hessian}
  \end{equation}
  Then the Hessian is $\nabla^2 L(\theta\opt_L, \sigma) = \Hessian(t_m)$.
\end{lemma}

\paragraph{Asymptotic normality of the minimizers.} 
%With the existence of minimizers assured, we turn to their asymptotics.
%For each of these, we require slightly different calculations, as the
%resulting covariances are differ.
To state asymptotic results when we have multiple labels,
we define the average link function $\wb{\sigma}\opt = \frac{1}{m} \sum_{j=1}^m \sigma_j\opt$ and the three functions
\begin{align}
	\linkerr(Z) & \defeq \sigma(t_m Z) (1 - \wb{\sigma}\opt(t\opt Z))
	- \sigma(-t_m Z) \wb{\sigma}\opt(t\opt Z), \nonumber \\
	\hessfunc(Z) & \defeq \sigma'(-t_m Z) \varphi_m(t\opt Z)
	+ \sigma'(t_m Z) (1 - \varphi_m(t\opt Z)),
	\label{eqn:error-functions}
	\\
	v_j(Z) & \defeq
	\sigma_j\opt(t\opt Z)(1 - \sigma_j\opt(t\opt Z))
	(\sigma(t_m Z) + \sigma(-t_m Z))^2.
	\nonumber
\end{align}
The first, the link error $\linkerr$, measures mis-specification of the
link $\sigma$ relative to the average $\wb{\sigma}\opt$. The second,
$\hessfunc$, is a Hessian term, as $\Hessian(t_m) =
\E[\hessfunc(Z) Z^2]u\opt{u\opt}^\top + \E[\hessfunc(Z)] \projperp \Sigma
\projperp$, and the third measures variance for labeler $j$.
We then prove the following theorem
in Appendix~\ref{sec:proof-multilabel-asymptotics}.
\begin{theorem}
  \label{theorem:multilabel-asymptotics}
  Let Assumptions~\ref{assumption:x-decomposition}
  and~\ref{assumption:model-link} hold, and let $\gme$ be the multi-label
  estimator~\eqref{eqn:general-estimators}.
  %
  Define the shorthand $\wb{v} =
  \frac{1}{m} \sum_{j = 1}^m v_j$. Then $\gme \cas \theta_L\opt$,
  and
  \begin{equation*}
    \sqrt{n} (\what{\theta}_n - \theta_L\opt)
    \cd \normal\left(0,
    \frac{\E[\linkerr(Z)^2 Z^2] + m^{-1} \E[\wb{v}(Z) Z^2]}{
      \E[\hessfunc(Z) Z^2]^2} u\opt {u\opt}^\top
    + \frac{\E[\linkerr(Z)^2] + m^{-1}\E[\wb{v}(Z)]}{\E[\hessfunc(Z)]^2}
    \prn{\projperp \Sigma \projperp}^\dagger
    \right).
  \end{equation*}
  Additionally, if $\what{u}_n = \gme / \ltwos{\gme}$
  and $t_m$ is the zero of the gap function $h_{t\opt, m}(t) = 0$, then
  \begin{equation*}
    \sqrt{n}(\what{u}_n - u\opt)
    \cd
    \normal\left(0, \frac{1}{t_m^2}
    \frac{\E[\linkerr(Z)^2] + m^{-1}
      \E[\wb{v}(Z)]}{\E[\hessfunc(Z)]^2}
    \left(\projperp \Sigma \projperp\right)^\dagger\right).
  \end{equation*}
\end{theorem}

Theorem~\ref{theorem:multilabel-asymptotics} exhibits two dependencies: the
first on the link error terms $\E[\linkerr(Z)^2]$---a bias
term---and the second on the rescaled average variance $\frac{1}{m}
\E[\wb{v}(Z)]$. The multi-label estimator recovers
an optimal $O(1/m)$ covariance if the link errors are negligible, but if
not, it necessarily has $O(1)$ asymptotic covariance.
%
%% The next corollary highlights how things simplify.
%
In the well-specified
case that $\sigma$ is symmetric and $\sigma = \wb{\sigma}\opt$, 
the zero of the gap
function~\eqref{eqn:h-m-func} is evidently $t_m = t\opt =
\ltwo{\theta\opt}$, the error term $\linkerr(Z) = 0$,
and $v_j(Z) = \sigma\opt(t\opt Z)(1 - \sigma\opt(t\opt Z))$, and by symmetry,
$\sigma'(t) = \sigma'(-t)$ so that
$\hessfunc(Z) = \sigma'(t\opt Z)$:
\begin{corollary}[The well-specified case] \label{cor:well-specified-normalized-cov}
	Let the conditions above hold. Then
	\begin{equation*}
		\sqrt{n}(\what{u}_n - u\opt)
		\cd \normal\left(0, \frac{1}{m} \cdot
		\frac{1}{\ltwo{\theta\opt}^2}
		\frac{\E[\sigma(t\opt Z) (1 - \sigma(t\opt Z))]}{\E[\sigma'(t\opt Z)]^2}
		\projperp \Sigma \projperp \right).
	\end{equation*}
\end{corollary}

\paragraph{Asymptotic normality with majority vote.}
When we use the majority vote estimators, the asymptotics differ:
no averaging reduces variance as $m$ increases, even in a
``well-specified'' case. The asymptotic variance typically decreases
$m$ grows, but at a slower rate roughly related to
the fraction of data where there is ``signal'', as we discuss
briefly following Corollary~\ref{cor:lowd-majority-vote-m}.
\begin{theorem}
  \label{theorem:majority-vote-asymptotics}
  Let Assumptions~\ref{assumption:x-decomposition}
  and~\ref{assumption:model-link} hold, and let $\what{\theta}_n = \mve$ be
  the general majority vote
  estimator~\eqref{eqn:general-estimators}. Let $t_m$ be
  the zero of the gap function~\eqref{eqn:h-m-func}, solving
  $h_{t\opt,m} (t) = 0$.  Then $\what{\theta}_n \cas
  \theta\opt_L = t_m u\opt$, and $\what{u}_n = \what{\theta}_n /
  \ltwos{\what{\theta}_n}$ satisfies
  \begin{equation*}
    \sqrt{n} \left(\what{u}_n - u\opt\right)
    \cd \normal\left(0, \frac{1}{t_m^2}
    \frac{\E[\sigma(-t_m |Z|)^2 \rho_m(t\opt Z) + \sigma(t_m |Z|)^2
	(1 - \rho_m(t\opt Z))]}{\E[\hessfunc(Z)]^2}
    \prn{\projperp \Sigma \projperp}^\dagger \right).
	\end{equation*}
\end{theorem}
\noindent
We defer the proof to Appendix~\ref{sec:proof-majority-vote-asymptotics}.

In most cases, we take a symmetric link function,
so that $\sigma(t) = 1 - \sigma(-t)$, and thus $\sigma'(t) = \sigma'(-t)$
and $\hessfunc(z) = \sigma'(t_m z) \ge 0$. This simplifies the denominator
in Theorem~\ref{theorem:majority-vote-asymptotics} to
$\E[\sigma'(t_m Z)]^2$.
Written differently, we may define a (scalar) variance-characterizing
function $\varfunc$ implicitly as follows: let $t_m = t_m(t)$ be a
zero of $h_{t,m}(s) =
\E[\sigma(sZ) Z (1 - \varphi_m(t Z))] - \E[\sigma(-s Z) Z \varphi_m(t Z)]$
in $s$, $h_{t,m}(t_m(t)) = 0$ so that $t_m$
is a function of the size $t$ (recall the gap~\eqref{eqn:h-m-func}),
and define
\begin{equation}
  \label{eqn:generic-majority-covariance}
  \varfunc(t) \defeq \frac{1}{t_m^2} \frac{\E[\sigma(-t_m |Z|)^2
      \rho_m(t Z) + \sigma(t_m |Z|)^2 (1 - \rho_m(t Z))]}{
    \E[\sigma'(t_m Z)]^2}
\end{equation}
where $t_m = t_m(t)$ above is implicitly defined.  Then
\begin{equation*}
  \sqrt{n}\left(\what{u}_n - u\opt\right)
  \cd \normal\left(0,
  \varfunc(t\opt) \prn{\projperp \Sigma \projperp}^\dagger
  \right).
\end{equation*}
Each of our main results, including those on well-specified models,
follows by characterizing the behavior of $\varfunc(t)$ in
the asymptotics as $m \to \infty$ and
the solution norm $t_m = \ltwos{\theta_L\opt}$ scales,
which the calibration gap~\eqref{eqn:h-m-func} determines.
The key is that the scaling with $m$ varies with
the fidelity of the model, behavior of the links $\sigma$, and
the noise exponent~\eqref{eqn:mammen-tsybakov}; our coming
consequences of the master Theorems~\ref{theorem:multilabel-asymptotics}
and~\ref{theorem:majority-vote-asymptotics} reify this scaling.

%% \jcdcomment{Highlight more here that we have to find scaling with $m$.}

%This may seem to imply that aggregated data perform better than non-aggregated method, and justify the widely usage of majority vote in practice at first glance. Indeed, our result says the majority-vote estimator $\mve$ is robust when the model is misspecified---however, as we will show in the next section, with a more careful designed approach, we can still have efficiency and perform better than majority vote-estimator as in the well-specified case.
%
%The message is the following: the majority-vote estimator (or other data aggregation method) seems tempting and easy to use, but both in the well-specified and misspecified case, we can do (and should do) better than the simple majority-vote estimator with non-aggregated data, calls for rethinking the general pipeline for dataset construction in modern machine learning.

\subsection{Robustness to model mis-specification}
\label{sec:robustness-model-misspecification}

Having established general convergence results for the multi-label
$\gme(\sigma)$ and majority vote  $\mve(\sigma)$ estimators, we 
explicate their performance when %% we have a mis-specified model---
the link $\sigma$ is incorrect. %% ---by leveraging
%% Theorems~\ref{theorem:multilabel-asymptotics}
%% and~\ref{theorem:majority-vote-asymptotics} to characterize their
%% asymptotics and show the robustness
%% of majority-vote estimators.
%% In particular, we consider our
%% logistic link based estimators of~\Cref{sec:well-sepcified-mod} ($\tmle$ and
%% $\mve$) which choose $\sigma = \sigma\lr$ regardless of the data and analyze
%% their performance. Ideally, we can try to estimate the true link functions
%% from the data---we defer this to~\Cref{sec:semi-param}.
%but we will focus on this simple setup here and discuss more in-depth approaches in the next section.
\paragraph{Multi-label estimator.}
To make interpretable statements about
the multi-label estimator $\gme$ in~\eqref{eqn:general-estimators},
assume identical links $\sigma_j\opt \equiv \sigma\opt \in
\flink$.
A corollary of Theorem~\ref{theorem:multilabel-asymptotics}
then follows:
\begin{corollary}
  \label{corollary:lowd-misspcfd-logistic-regression}
  Let Assumptions~\ref{assumption:x-decomposition} and
  \ref{assumption:z-density} hold for some $\beta \in (0, \infty)$, and
  $t\opt = \ltwos{\theta\opt}$. Then the calibration
  gap function~\eqref{eqn:h-m-func} has unique positive zero
  $t_{\sigma\opt}$ such that
  $h_{t\opt,m}(t_{\sigma\opt}) = 0$, and the multi-label
  estimator~\eqref{eqn:general-estimators} satisfies
  \begin{equation*}
    \gme \cp  t_{\sigma\opt} u\opt.
  \end{equation*}
  For the mean $\wb{v} = \frac{1}{m} \sum_{j = 1}^m v_j$ of the
  variances~\eqref{eqn:error-functions}), the estimate
  $\what{u}_{n,m} = \gme / \ltwos{\gme}$ satisfies
  \begin{equation*}
    \sqrt{n} \left(\what{u}_{n,m} - u\opt\right)
    \cd \normal\left(0,
    \frac{\E[\linkerr(Z)^2]  + m^{-1} \E[\wb{v}(Z)]}{t_{\sigma\opt}^2
      \E[\hessfunc(Z)]^2}  
    % C_{m,\sigma\opt}(t_{\sigma\opt})
    \prn{\proj_{u\opt} \Sigma \proj_{u\opt}}^\dagger\right)  .
  \end{equation*}
  %% satisfies $\lim_{t
  %%   \to \infty} \lim_{m \to \infty} C_{m,\sigma\opt}(t) t^{2-2\beta}
  %% = c_Z \E[|Z|^{\beta - 1} \sigma'(|Z|)]$ and
\end{corollary}

In this simplified case, the asymptotic covariance remains of constant
order in $m$ unless $\E[\linkerr(Z)^2] = 0$.  In contrast,
the majority vote estimator exhibits more robustness; this is perhaps
expected, as Corollary~\ref{cor:lowd-majority-vote-m} shows covariance scaling as $1/\sqrt{m}$ in the
logistic link case, which is \emph{a fortiori} misspecified for majority
vote labels.
%has covariance scaling as $1/\sqrt{m}$, though the
%generality of the behavior and its distinction from
%Corollary~\ref{corollary:lowd-misspcfd-logistic-regression}
%is interesting.

\paragraph{Majority vote estimator.}
For the majority-vote estimator, we relax our assumptions and allow
$\sigma_j\opt$ to differ, showing how the broad conclusions
Corollary~\ref{corollary:lowd-misspcfd-logistic-regression} suggests
hold in some generality: majority vote estimators achieve slower
convergence than well-specified (maximum likelihood) estimators using each
label but exhibit more robustness.
%
To characterize the large $m$
behavior, we require the following regularity conditions on
the average link $\wb{\sigma}_m\opt = \frac{1}{m} \sum_{j=1}^m \sigma_j\opt$.
%
%% which we require has a limiting derivative at $0$.

\begin{assumption} \label{assumption:condition-misspcfd-link-majority-vote}
  For the sequence of link functions $\{\sigma_j \mid j \in \mathbb{N}\}
  \subset \flink$, let $\wb{\sigma}_m\opt = \frac{1}{m} \sum_{j=1}^m
  \sigma_j\opt$. There exists ${\wb{\sigma}\opt}'(0) >0$ such that
  \begin{subequations}
    \begin{align}
      \bm{\mathrm{(i)}} \qquad  & \lim_{m \to \infty} \sqrt{m} \prn{\wb{\sigma}_m\opt\prn{\frac{t}{\sqrt m}} - \frac 1 2} = {\wb{\sigma}\opt}'(0) t   , \quad \textrm{for each } t \in \R  ; \label{ass:misspcfd-link-1} \\
      \bm{\mathrm{(ii)}} \qquad &\liminf_{m \to \infty} \inf_{t \neq 0} \frac{\left|\wb{\sigma}_m\opt(t) - \frac 1 2 \right|}{\min \left\{|t|, 1\right\}} > 0  ; \label{ass:misspcfd-link-2} \\
      \bm{\mathrm{(iii)}} \qquad &\lim_{t \to 0} \sup_{j \in \N}
      \left|\sigma_j\opt(t) - \frac 1 2 \right| = 0  . \label{ass:misspcfd-link-3}
    \end{align}
  \end{subequations}
\end{assumption}
\noindent
These assumptions simplify if the links are identical:
if $\sigma_j\opt \equiv \sigma\opt$,
we only require $\sigma\opt$ is differentiable around
$0$ with ${\sigma\opt}'(0) >0$ and $|\sigma\opt(t) - \half|
\gtrsim \min\{|t|, 1\}$.

%\ha{I think it might be worth it to only consider this simpler setting where all $\sigma_j$ are identical and postpone the general case to appendix. This will remove Assumption A4 and simplify things a bit.}

We can apply Theorem~\ref{theorem:majority-vote-asymptotics} to obtain
asymptotic normality for the
estimator~\eqref{eqn:general-estimators}. Recall
\begin{equation}
  \label{eq:link-function-m-majority-vote-general}
  \rho_m(t) \defeq
  \Prb(\majY = \sgn(\<X, \theta^\star\>) \mid \<X,
  \theta\opt\> = t),
\end{equation}
the probability of a correct majority vote given the margin $\<X,
\theta\opt\> = t$, and the calibration gap function~\eqref{eqn:h-m-func},
which by a calculation resolves to the more convenient form
\begin{align*}
  h(t) & = h_{t\opt, m}(t) = \E[\sigma(t|Z|) |Z| (1 - \rho_m(t\opt Z))]
  - \E[\sigma(-t |Z|) |Z| \rho_m(t\opt Z)].
\end{align*}
The main challenge is to characterize the large $m$ behavior for
the asymptotic covariance function $\varfunc(t)$ the
quantity~\eqref{eqn:generic-majority-covariance} implicitly
defines.
%
We postpone details to Appendix~\ref{proof:lowd-misspcfd-majority-vote-m}
and state the result, which follows from
Theorem~\ref{theorem:majority-vote-asymptotics} and an asymptotic expansion
of $\varfunc(t)$.

\begin{proposition}
  \label{prop:lowd-misspcfd-majority-vote-m}
  Let Assumptions~\ref{assumption:x-decomposition} and
  \ref{assumption:z-density} hold for some $\beta \in (0, \infty)$ with
  $\int_0^\infty z^{\beta - 1} \sigma'(z) dz < \infty$
  and
  $t\opt = \ltwos{\theta\opt}$, and in addition that
  Assumption~\ref{assumption:condition-misspcfd-link-majority-vote} holds
  and $\sigma$ is symmetric.
  Then there are constants $a, b > 0$, depending only on $\beta$, $c_Z$,
  $\sigma$, and ${\wb{\sigma}\opt}'(0)$, such that there is a unique $t_m
  \geq t_1 = t^\star$ solving $h(t_m) = 0$, and for this $t_m$ we have both
  \begin{equation*}
    \what{\theta}\mv_{n,m} \cp t_m u\opt
    ~~ \mbox{and} ~~
    \lim_{m \to \infty} \frac{t_m}{t\opt \sqrt{m}} = a  .
  \end{equation*}
  Moreover, the covariance~\eqref{eqn:generic-majority-covariance} has the
  form $\varfunc(t) = \frac{b}{(t \sqrt m)^{2 - \beta}}(1 + o_m(1))$, and
  \begin{align*}
    \sqrt{n} \left(\what{u}\mv_{n,m} - u^\star \right)
    &\cd \normal\left(0, C_m(t\opt) \prn{\proj_{u\opt} \Sigma \proj_{u\opt}}^\dagger\right).
  \end{align*}
\end{proposition}

Proposition~\ref{prop:lowd-misspcfd-majority-vote-m} highlights the
majority vote estimator's robustness: even when the link $\sigma$ is more
or less arbitrarily incorrect, the asymptotic covariance
scales reasonably. The noise parameter $\beta$ in
Assumption~\ref{assumption:z-density}, analogous
to the noise exponent~\eqref{eqn:mammen-tsybakov}, plays
an important role. In typical cases with $\beta = 1$ (e.g., when
$X \sim \normal(0, I_d)$), we see $\varfunc(t) \asymp \frac{1}{t \sqrt{m}}$.
In noisier cases, corresponding to $\beta \downarrow 0$, majority vote
provides benefits approaching a well-specified model; conversely,
in ``easy'' cases where $\beta > 2$, majority vote estimators become
\emph{more} unstable, as they make the data nearly separable,
making margin-type estimators
volatile~\cite{CandesSu19pnas}.
